\subsection{单位元}

定义1. 在一个合成律 \(\top\) （定义在集合 \(\mathbf{E}\) 上）下，一个元素 \(e\) （属于 \(\mathbf{E}\) ）称为单位元，如果对于所有 \(x \in \mathbf{E}\) , \(e \top x = x \top e = x\) .

在给定的律 \(\top\) 下至多存在一个单位元，因为如果 \(e\) 和 \(e'\) 都是单位元，则 \(e = e \top e' = e'\) 。单位元与每个元素都可交换：它是一个中心元素。

定义2. 具有单位元的广群称为含幺广群。若 E、E' 是含幺广群，把广群 E 映入广群 E' 且将 E 的单位元映到 E' 的单位元的同态称为 E 到 E' 的保单位同态（或态射）。结合的含幺广群称为幺半群。

若E、E'是幺半群，则从E到E'的保单位态射称为幺半群同态或幺半群态射。

示例。(1) 在自然数的集合 \(\mathbf{N}\) 中，0 是加法的单位元，1 是乘法的单位元。这两条律各自给出 \(\mathbf{N}\) 一个交换幺半群结构（集合论，III，§3，第3号）。

\begin{enumerate}
\def\labelenumi{(\arabic{enumi})}
\setcounter{enumi}{1}
\item
  在集合 E 的子集族中， \(\varnothing\) 在律 \(\cup\) 下是单位元，而 E 在律 \(\cap\) 下是单位元。更一般地，在一个格中，最小元素（若存在）在上确界的律下是单位元；反之，若在该律下存在单位元，则它是集合的最小元素。对最大元素与下确界的律亦有类似结论。
\item
  集合 \(\mathbf{N}\) 在律 \((x, y) \mapsto x^y\) 下没有单位元。在律 \((\mathbf{X}, \mathbf{Y}) \mapsto \mathbf{X} \circ \mathbf{Y}\) （在 \(\mathbf{E} \times \mathbf{E}\) 的子集之间）下，对角映射 \(\Delta\) 是单位元。在律 \((f, g) \mapsto f \circ g\) （在 \(\mathbf{E}\) 到 \(\mathbf{E}\) 的映射之间）下， \(\mathbf{E}\) 到 \(\mathbf{E}\) 上的恒等映射是单位元。
\item
  设 E 是一个广群，R 是 E 上与 E 上的律相容的一个等价关系（§ 1，第 6 号）。若 e 是 E 的一个单位元，则 e 在 E/R 中的典范像就是广群 E/R 的一个单位元。
\end{enumerate}

含幺广群的单位元是一个保单位同态；两个保单位同态的复合也是一个保单位同态。一个映射成为含幺广群同构的充要条件是它是一个双射的保单位同态，并且此时逆映射也是一个保单位同态。设 E 和 E' 为含幺广群， \(e'\) 为 \(E'\) 的单位元；将 E 映到 \(E'\) 的常值映射（它把 E 映到 \(e'\) ）是一个保单位同态，称为平凡同态。

一族含幺广群（相应地，幺半群）的积是含幺广群（相应地，幺半群）。

每个含幺广群（相应地，幺半群）的商广群都是含幺广群（相应地，幺半群）。

设 E 是一个含幺广群，e 为其单位元。E 的一个子广群 A 若满足 \(e \in A\) 则称为 E 的含幺子广群。显然，e 是广群 A 的单位元。E 的含幺子广群的任意交集仍是 E 的含幺子广群。若 X 是 E 的子集，则存在包含 X 的 E 的最小含幺子广群；它称为由 X 生成的 E 的含幺子广群；当 X 为空时它等于 \(\{e\}\) 。若 E 是一个幺半群，则 E 的含幺子广群称为 E 的子幺半群。

如果 F 是一个没有单位元的广群，F 的子广群可能具有单位元。例如，如果 F 是结合的，且 h 是 F 的幂等元（I，§ 1，第 4 号），则集合 \(h \top x \top h\) ，其中 x 遍历 F，是 F 的一个以 h 为单位元的子广群。

若 E 是带单位元 e 的广群，则 E 的一个子广群 A 即使满足 \(e \notin A\) 也可能仍具有单位元。

定义3. 设 E 为一个含幺广群。E 的单位元称为 E 的元素的空族的合成。

若 \((x_{\alpha})_{\alpha \in \varnothing}\) 是 E 的元素组成的空族，其合成 \(e\) 也记作 \(\bigcap_{\alpha \in \varnothing} x_{\alpha}\) 。例如，我们写

\[
\bigcap_ {q \leqslant i \leqslant p} x _ {i} = e
\]

当 \(p < q\) ( \(p, q \in \mathbf{N}\) ) 时。类似地，我们写 \(\top^0 x = e\) （对任意的 \(x\) ）。有了这些定义，即使略去集合 \(\mathbf{A}\) 和 \(\mathbf{B}_i\) 非空这一假设，§1 中的定理1和定理3仍然成立。类似地，公式 \(\top^{m+n} x = \left( \frac{m}{\top} x \right)^{\top} \left( \frac{n}{\top} x \right)\) 和 \(\top^{mn} x = \frac{m}{\top} \left( \frac{n}{\top} x \right)\) 对 \(m \geqslant 0, n \geqslant 0\) 成立.

设 E 是一个含幺广群，其律记为 \(\top\) ， \(e\) 为其单位元。E 的元素的一族 \((x_{i})_{i\in I}\) 的支集是指标 \(i\in I\) 中使 \(x_{i}\neq e\) 者所成的集合。设 \((x_{i})_{i\in I}\) 为 E 中具有有限支集的一族元素。我们将在以下两种情形中定义合成 \(\bigcap_{i\in I}x_{i}\) ：

\begin{enumerate}
\def\labelenumi{(\alph{enumi})}
\item
  集合 I 是全序的；
\item
  E 是结合的，且 \(x_{i}\) 是两两可交换的。
\end{enumerate}

在这两种情形下，令 S 为族 \((x_{i})\) 的支集。若 J 是 I 的包含 S 的有限子集，则 \(\bigcap_{i\in J}x_{i} = \bigcap_{i\in S}x_{i}\) ，这可通过对 J 中元素个数进行归纳看出：在情形 (a) 中应用 §1 的定理 1，在情形 (b) 中应用 §1 的定理 3。设 \(\bigcap_{i\in I}x_{i}\) 表示所有包含 S 的 I 的有限子集 J 的合成 \(\bigcap_{i\in J}x_{i}\) 的公共值。当 I 是区间 \(\{p,\rightarrow\{\text{的 }N\}\) 时，我们也记作 \(\bigcap_{i=p}^{\infty}x_{i}\) .

有了这些定义和记号，§1中的定理1和定理3以及定理3后的注记可以推广到具有有限支集的族。

以加法记法书写的律的单位元通常记作 0，称为零元或零元素（有时也称为原点）。在以乘法记法书写的律下，它通常记作 1，称为单位元（或单位）。

